% Partie : logiques multivaluées
% Chapitre : calcul des séquents
% Section : notions de théorie de la démonstration

\documentclass[../../../include/open-logic-section]{subfiles}

\begin{document}

\olfileid[fr]{mvl}{seq}{prf}

\olsection{Notions de théorie de la démonstration}

Cette section porte sur le calcul classique $\Log{LK}$ à deux
côtés, avec un antécédent et un succédent.








\begin{explain}
Après avoir défini plusieurs notions sémantiques importantes
— validité, conséquence logique et satisfaisabilité —,
nous définissons les \emph{notions de théorie de la
démonstration} correspondantes. Celles-ci font appel à la
!!{derivability} ou à la !!{nonderivability} de certains
séquents, plutôt qu'à la satisfaction d'!!{sentence}s dans
des !!{structure}s. La coïncidence de ces notions a constitué
une découverte importante. Elle est exprimée par les
théorèmes de \emph{correction} et de \emph{complétude}.
\end{explain}

\begin{defn}[Théorèmes]
!!^a{sentence}~$!A$ est un \emph{théorème} s'il existe
!!a{derivation} dans~$\Log{LK}$ du séquent
$\quad\Sequent !A$.
Nous écrivons $\Proves !A$ si $!A$ est un théorème et
$\Proves/ !A$ sinon.
\end{defn}

\begin{defn}[!!^{derivability}]
!!^a{sentence}~$!A$ est \emph{!!{derivable} à partir d'un}
ensemble d'!!{sentence}s~$\Gamma$, ce que l'on note
$\Gamma\Proves !A$, si et seulement s'il existe un
sous-ensemble fini~$\Gamma_0\subseteq\Gamma$ et une suite
$\Gamma_0'$ des !!{sentence}s de~$\Gamma_0$ tels que
$\Log{LK}$ !!{derive}s le séquent $\Gamma_0'\Sequent !A$.
Si $!A$ n'est pas !!{derivable} à partir de $\Gamma$,
nous écrivons $\Gamma\Proves/ !A$.
\end{defn}

Les règles de contraction, d'affaiblissement et d'échange
rendent indifférents l'ordre et le nombre d'occurrences des
!!{sentence}s dans~$\Gamma_0'$.
Si $\Gamma_0'\Sequent !A$ est !!{derivable}, alors
$\Gamma_0''\Sequent !A$ l'est aussi pour toute suite
$\Gamma_0''$ contenant les mêmes !!{sentence}s que
$\Gamma_0'$.
Par exemple, si $\Gamma_0=\{!B,!C\}$, les suites
$\Gamma_0'=\tuple{!B,!B,!C}$ et
$\Gamma_0''=\tuple{!C,!C,!B}$ contiennent exactement les
!!{sentence}s de~$\Gamma_0$. Si un séquent ayant l'une pour
antécédent est !!{derivable}, celui ayant l'autre pour
antécédent l'est également :
\begin{prooftree}
  \AxiomC{}
  \Deduce$!B, !B, !C \fCenter !A$
  \RightLabel{\LeftR{\Contraction}}
  \UnaryInf$!B, !C \fCenter !A$
  \RightLabel{\LeftR{\Exchange}}
  \UnaryInf$!C, !B \fCenter !A$
  \RightLabel{\LeftR{\Weakening}}
  \UnaryInf$!C, !C, !B \fCenter !A$
\end{prooftree}
Désormais, lorsque $\Gamma_0$ est un ensemble fini d'%
!!{sentence}s, $\Gamma_0\Sequent !A$ désigne un séquent dont
l'antécédent est une suite dont les termes sont exactement
les !!{sentence}s de~$\Gamma_0$, avec répétitions éventuelles.





Les contractions, échanges et affaiblissements nécessaires
seront employés tacitement.

\begin{defn}[Cohérence]
Un ensemble d'!!{sentence}s~$\Gamma$ est \emph{incohérent} si
et seulement s'il existe un sous-ensemble fini
$\Gamma_0\subseteq\Gamma$ tel que $\Log{LK}$ !!{derive}s
$\Gamma_0\Sequent\quad$.
S'il n'est pas incohérent, c'est-à-dire si, pour tout
sous-ensemble fini $\Gamma_0\subseteq\Gamma$,
$\Log{LK}$ ne !!{derive}s pas $\Gamma_0\Sequent\quad$,
nous disons qu'il est \emph{cohérent}.
\end{defn}

\begin{prop}[Réflexivité]
\ollabel{prop:reflexivity}
Si $!A\in\Gamma$, alors $\Gamma\Proves !A$.
\end{prop}

\begin{proof}
Le séquent initial $!A\Sequent !A$ est !!{derivable}, et
$\{!A\}\subseteq\Gamma$.
\end{proof}

\begin{prop}[Monotonie]
\ollabel{prop:monotonicity}
Si $\Gamma\subseteq\Delta$ et $\Gamma\Proves !A$, alors
$\Delta\Proves !A$.
\end{prop}

\begin{proof}
Supposons que $\Gamma\Proves !A$. Il existe donc un ensemble
fini $\Gamma_0\subseteq\Gamma$ tel que
$\Gamma_0\Sequent !A$ soit !!{derivable}.
Puisque $\Gamma\subseteq\Delta$, $\Gamma_0$ est aussi un
sous-ensemble fini de~$\Delta$.
La !!{derivation} de $\Gamma_0\Sequent !A$ établit donc
également que $\Delta\Proves !A$.
\end{proof}

\begin{prop}[Transitivité]
\ollabel{prop:transitivity}
Si $\Gamma\Proves !A$ et $\{!A\}\cup\Delta\Proves !B$,
alors $\Gamma\cup\Delta\Proves !B$.
\end{prop}

\begin{proof}
Si $\Gamma\Proves !A$, il existe un sous-ensemble fini
$\Gamma_0\subseteq\Gamma$ et !!a{derivation}~$\pi_0$
de $\Gamma_0\Sequent !A$.
Si $\{!A\}\cup\Delta\Proves !B$, il existe un sous-ensemble
fini $\Delta_0\subseteq\Delta$ et !!a{derivation}~$\pi_1$
de $!A,\Delta_0\Sequent !B$, en ajoutant au besoin la
formule de coupure par affaiblissement.



Considérons la !!{derivation} suivante :
\begin{prooftree}
  \AxiomC{}
  \RightLabel{$\pi_0$}
  \Deduce$\Gamma_0 \fCenter !A$
  \AxiomC{}
  \RightLabel{$\pi_1$}
  \Deduce$!A, \Delta_0 \fCenter !B$
  \RightLabel{\Cut}
  \BinaryInf$\Gamma_0, \Delta_0 \fCenter !B$
\end{prooftree}
Puisque $\Gamma_0\cup\Delta_0\subseteq\Gamma\cup\Delta$,
elle établit $\Gamma\cup\Delta\Proves !B$.
\end{proof}

En particulier, si $\Gamma\Proves !A$ et $!A\Proves !B$,
alors $\Gamma\Proves !B$.
Il en découle aussi que, si $!A_1,\dots,!A_n\Proves !B$
et $\Gamma\Proves !A_i$ pour chaque~$i$, alors
$\Gamma\Proves !B$.

\begin{prop}
\ollabel{prop:incons}
$\Gamma$ est incohérent si et seulement si
$\Gamma\Proves {!A}$ pour tout !!{sentence}~$!A$.
\end{prop}

\begin{proof}
Exercice.
\end{proof}

\begin{prob}
Démontrer la \olref[mvl][seq][prf]{prop:incons}.\footnote{Correction
de l'édition française : la référence de la source renvoie
à une autre section. Elle vise ici la proposition immédiatement
précédente.}
\end{prob}

\begin{prop}[Compacité]
\ollabel{prop:proves-compact}
\begin{enumerate}
\item Si $\Gamma\Proves !A$, il existe un sous-ensemble fini
$\Gamma_0\subseteq\Gamma$ tel que $\Gamma_0\Proves !A$.
\item Si tout sous-ensemble fini de~$\Gamma$ est cohérent,
alors $\Gamma$ est cohérent.
\end{enumerate}
\end{prop}

\begin{proof}
\begin{enumerate}
\item Si $\Gamma\Proves !A$, il existe un sous-ensemble fini
$\Gamma_0\subseteq\Gamma$ tel que le séquent
$\Gamma_0\Sequent !A$ ait !!a{derivation}.
Par conséquent, $\Gamma_0\Proves !A$.
\item Si $\Gamma$ est incohérent, il existe un sous-ensemble
fini~$\Gamma_0\subseteq\Gamma$ tel que $\Log{LK}$ !!{derive}s
$\Gamma_0\Sequent\quad$. Mais $\Gamma_0$ est alors un
sous-ensemble fini de~$\Gamma$ qui est incohérent.
\end{enumerate}
\end{proof}

\end{document}
